\documentclass[11pt]{article}
\usepackage[margin=1in]{geometry}
\usepackage{amsmath,amssymb,amsthm,mathtools}
\usepackage{empheq}
\usepackage{booktabs,multirow}
\usepackage{enumitem}
\usepackage{graphicx}
\usepackage{xcolor}
\usepackage[T1]{fontenc}
\usepackage{titlesec}
\usepackage{authblk}
\usepackage[numbers,sort&compress]{natbib}
\usepackage{hyperref}
\hypersetup{colorlinks=true, linkcolor=blue!45!black, citecolor=blue!45!black,
            urlcolor=blue!45!black}
\allowdisplaybreaks

% ---- SIAM-style run-in section headings (Charous & Lermusiaux 2023) ---- %
\titleformat{\section}[runin]{\normalfont\bfseries}{\thesection.}{0.5em}{}[.\quad]
\titleformat{\subsection}[runin]{\normalfont\bfseries}{\thesubsection.}{0.5em}{}[.\quad]
\titleformat{\subsubsection}[runin]{\normalfont\itshape}{\thesubsubsection.}{0.5em}{}[.\quad]
\titlespacing*{\section}{0pt}{1.4em}{0.4em}
\titlespacing*{\subsection}{0pt}{1.0em}{0.4em}
\titlespacing*{\subsubsection}{0pt}{0.7em}{0.4em}

\theoremstyle{plain}
\newtheorem{proposition}{Proposition}
\newtheorem{lemma}{Lemma}
\newtheorem{corollary}{Corollary}
\theoremstyle{definition}
\newtheorem{remark}{Remark}
\newtheorem{assumption}{Assumption}

\newcommand{\TODO}[1]{\textcolor{red}{\textbf{[TODO: #1]}}}
\newcommand{\figph}[2]{%
  \begin{center}\setlength{\fboxsep}{10pt}%
  \fcolorbox{gray!60}{gray!8}{\parbox[c][#1][c]{0.82\linewidth}{%
  \centering\itshape\color{gray!50!black} #2}}\end{center}}

\newcommand{\Lop}{\mathcal{L}}
\newcommand{\Jop}{\mathcal{J}}
\newcommand{\Ndot}{\dot{N}}
\newcommand{\Nddot}{\ddot{N}}
\newcommand{\Ndddot}{\dddot{N}}
\newcommand{\bO}{\mathcal{O}}
\newcommand{\DT}{\Delta T}
\newcommand{\DTs}{\Delta T^{\star}}
\newcommand{\hf}{h_{\mathrm{f}}}
\newcommand{\taul}{\tau_{\lambda}}
\newcommand{\norm}[1]{\left\lVert #1 \right\rVert}
\newcommand{\E}{\mathbb{E}}
\newcommand{\zL}{z_L}
\newcommand{\zN}{z_N}
\newcommand{\shat}{\hat{\sigma}}

\title{\bfseries A Step-Size Resolution Closure for IMEX Integration of\\
Two-Dimensional Quasi-Geostrophic Turbulence:\\
\large Truncation-Error Analysis, Wiener-Conditioned Learned Stencils,\\
and Von Neumann--Certified Stability}

\author[1]{Sanaa Mouzahir}
\author[1]{\TODO{co-authors / advisor}}
\affil[1]{\small Department of Mechanical Engineering and Center for
Computational Science and Engineering,\\ Massachusetts Institute of Technology,
Cambridge, MA 02139}
\date{\today}

\begin{document}
\maketitle

\begin{abstract}
\noindent
We treat the \emph{time}-discretization error of a fixed numerical scheme as a
closure problem. For the spectral Adams--Bashforth-2 / Crank--Nicolson
(AB2--CN2) IMEX integrator applied to two-dimensional quasi-geostrophic (QG)
turbulence, we derive the local truncation error in closed form to
$\bO(\DT^6)$ and compare it, term by term, with those of AB4--CN2 and
classical RK4. Each truncation operator splits into an analytical part (powers
of the Fourier-diagonal linear operator, free at inference) and an expensive
part built from time derivatives of the advective Jacobian; a compact
($\approx$3{,}700-parameter) physics-initialized network estimates only the
latter from a seven-level time stencil. The estimator's error is analyzed
exactly: the finite-difference-in-time error at every order is driven by the
single field $\omega^{(S)}$, and a Wiener (least-squares) analysis shows the
optimal learned stencil correction varies across step sizes and flow regimes as
$(\DT\,\sigma(\kappa))^{S-k}$, where $\sigma(\kappa)$ is a per-shell
decorrelation rate measurable from the model's own inputs to
$\bO(10^{-3})$ relative accuracy. This yields a conditioned architecture whose
pooled training floor deletes the dominant (pooled-variance) term of the
unconditioned floor. Stability is handled symmetrically: the modified-equation
analysis shows that a closure targeting the exact flow has zero linear
stability margin on the advection axis, while targeting fine-step RK4 buys a
provable $y^6/144$ margin; a differentiable frozen-coefficient von Neumann
certificate of the \emph{closed} scheme is derived and used both as an a
priori diagnostic and as a training penalty. The analytic (network-free) $R_3$
closure is unconditionally stable in our experiments and improves coarse-step
accuracy by one to two orders of magnitude; the learned closure attains a
replicated median
$16\times$ error reduction over ten initial conditions at $\DT=5\times10^{-3}$
on a 16-step horizon, while a horizon study shows the gain is not yet
unconditional --- blow-ups reappear at $\sim$2$\times$ the fine-tuned horizon
--- which we diagnose (a network-noise history-contamination loop) and
attack with the certificate penalty. We report the complete error, cost, and
stability budget of the method, including its failure modes.

\smallskip
\noindent\textbf{Key words.} time-integration error, modified equation,
closure modeling, IMEX schemes, Wiener filtering, von Neumann stability,
pseudo-spectral methods, quasi-geostrophic turbulence, scientific machine
learning
\end{abstract}

% ===================================================================== %
\section{Introduction}
\label{sec:intro}
Numerical simulations of geophysical turbulence are limited by the time step.
In stiff pseudo-spectral solvers the per-step cost is dominated by the
nonlinear (advective) term, and the total cost scales as $T/\Delta t$ to reach
a horizon $T$. For an integrator whose accuracy---not linear stability---sets
$\Delta t$, halving the usable step doubles the cost. Reducing this
time-discretization cost \emph{without changing the discretization itself} is
the problem addressed here.

Data-driven closures are by now standard for unresolved \emph{spatial} scales
on coarse grids, in engineering large-eddy simulation and geophysical modeling
alike \citep{gupta2021,guan2022,srinivasan2024}. The analogous problem for
unresolved \emph{temporal} scales---the error incurred by taking a large time
step with a fixed scheme---has received far less attention. We treat it as a
closure problem: derive the scheme's local truncation error (LTE) in closed
form, identify the single piece that is expensive to evaluate, and learn only
that piece.

\paragraph{Goal.}
The target of the construction is stated once and drives every design choice
in the paper:
\begin{quote}\itshape
run at the per-step cost of AB2--CN2 with a coarse step $\DT$, while
recovering the accuracy---and, deliberately, the stability margin---of
fine-step RK4 at $\hf=\DT/K$.
\end{quote}
The scheme and its formal order are unchanged; a learned correction makes one
coarse step behave as if the same trajectory had been advanced by $K$ fine RK4
steps. Because a fine-step RK4 trajectory coincides with the exact flow to
near machine precision, the $K\to\infty$ limit targets the exact flow;
Section~\ref{sec:closure} shows why the RK4 target, not the exact flow, is the
right deployment choice: on the advection axis the exact flow has zero
stability margin, while RK4's faint dissipation buys a provable $y^6/144$
buffer against both the analytic truncation tail and the network error.

\paragraph{Relation to classical ideas.}
The construction sits between three classical methods while remaining distinct
from each. Modified-equation analysis \citep{warming1974} supplies the
truncation machinery but is normally used to characterize or add order, not to
emulate a finer-step version of the \emph{same} scheme. Defect correction
\citep{stetter1978} iteratively matches a higher-order scheme at the same
step; we match the same scheme (or RK4) at a smaller step, a target with no
defect-correction analogue---the gain is set by $K$ at training time, not
bounded by a reference order. Parareal \citep{lions2001} attains
fine-propagator accuracy at coarse cost by iterating over the horizon; our
closure is an \emph{amortized} parareal in which the fine propagator is paid
for once, during training, and reused at inference without iteration.
High-order time integration that pays once for the curvature of a constrained
problem was developed for dynamical low-rank systems in the same group
\citep{charous2023}. On the learning side, the estimator belongs to the
learned-discretization lineage \citep{barsinai2019,kochkov2021}: a small
network constrained to the algebraic structure of the operator it replaces,
not a generic operator learner.

\paragraph{Contributions.}
(i) Closed-form LTEs of AB2--CN2, AB4--CN2 and RK4 for the QG vorticity
equation, in a shared elementary-differential language, verified symbolically
and by convergence studies (Sec.~\ref{sec:error}).
(ii) A quantitative validity theory for any such closure: the truncation
series has a finite convergence radius $\DTs$, measured per regime and
collapsed to a universal law $\DTs = C\,\taul$, $C=2.08\pm0.15$, with $\taul$
the Taylor microscale of the nonlinear term; a finite-lag stencil adds an
inner wall at span $\approx\DTs$; and error propagation through the closure is
benign (1:1 in the $\Nddot$ error) across the whole trained envelope
(Sec.~\ref{sec:error-magnitudes}).
(iii) A cost and stability budget of the schemes and their closed variants,
including exact amplification factors for the IMEX family and the margin
asymmetry between the exact-flow and RK4 targets
(Secs.~\ref{sec:cost}--\ref{sec:stability}, \ref{sec:closure}).
(iv) A structure-matched estimator of the expensive LTE part
($3{,}700$ parameters, exact-physics initialization) together with an
\emph{exact} error analysis of the estimator itself, culminating in a Wiener
floor decomposition that motivates---and quantifies the gain of---physics
conditioning of the learned stencils on a per-sample measurable rate
$\shat(\kappa)$ (Sec.~\ref{sec:method}).
(v) A differentiable frozen-coefficient von Neumann certificate for the
\emph{closed} scheme, used as a training penalty (Sec.~\ref{sec:nn-stability}).
(vi) An honest experimental record: replicated $16\times$ coarse-step error
reduction at $\DT=5\times10^{-3}$ over 16-step horizons, an analytic-closure
ceiling of $19$--$130\times$, and the two failure modes that remain---the
large-$\DT$ stability gap and the horizon wall---with their diagnosed
mechanism (Sec.~\ref{sec:results}).

\paragraph{Notation.}
Table~\ref{tab:notation} collects recurring symbols.
\begin{table}[h]
\centering\small
\caption{Key notation.}
\label{tab:notation}
\begin{tabular}{ll}
\toprule
symbol & meaning \\
\midrule
$\omega,\ \psi$        & vorticity and streamfunction, $\psi=\Delta^{-1}\omega$ \\
$\Lop$                 & stiff linear operator $\nu\Delta-\mu-\beta\partial_x\Delta^{-1}$ (Fourier-diagonal) \\
$N(\omega)$            & nonlinear term $-\Jop(\psi,\omega)+F$, $F$ time-independent \\
$\Jop(a,b)$            & advective Jacobian $\partial_xa\,\partial_yb-\partial_ya\,\partial_xb$ \\
$B(a,b)$               & symmetric bilinear form with $B(\omega,\omega)=-\Jop(\Delta^{-1}\omega,\omega)$ \\
$Q^{(k)},\ \dot Q$     & $k$-th total time derivative of $Q$ along the trajectory \\
$\DT,\ \hf,\ K$        & coarse step, fine step, ratio $K=\DT/\hf$; $h$ a generic step \\
$S$                    & time-stencil depth (number of stored levels; $S=7$ in production) \\
$\tau$                 & per-step local truncation error; $R_p$, $D_p$ its order-$p$ operator and denominator \\
$\DTs,\ \taul$         & convergence radius of the truncation series; Taylor microscale of $N$ \\
$\sigma(\kappa)$       & per-shell decorrelation rate $\norm{\dot\omega_\kappa}/\norm{\omega_\kappa}$; $\shat$ its 2-mark estimate \\
$\zL,\zN,z$            & per-mode eigenvalues $h\lambda_L$, $h\lambda_N$, $z=\zL+\zN$ \\
$g(z)$                 & amplification factor; $R(z)$ the RK4 stability polynomial \\
$f_{\rm NN}$           & learned (expensive) part of the closure \\
\bottomrule
\end{tabular}
\end{table}

\paragraph{Outline.}
Section~\ref{sec:equations} states the governing equations and the exact
derivative recursions every later formula is built from.
Section~\ref{sec:schemes} defines the three schemes.
Section~\ref{sec:scheme-analysis} analyzes them: error
(\ref{sec:error}--\ref{sec:error-magnitudes}), cost (\ref{sec:cost}), and
stability (\ref{sec:stability}). Section~\ref{sec:closure} assembles the
step-size resolution closure and its two possible targets.
Section~\ref{sec:method} presents the learned estimator: architecture
(\ref{sec:arch}), exact error analysis and Wiener conditioning
(\ref{sec:wiener}--\ref{sec:conditioning}), and the loss including the von
Neumann penalty (\ref{sec:loss}). Section~\ref{sec:nn-stability} analyzes the
stability of the closed scheme. Section~\ref{sec:training} documents the
ensemble training pipeline. Section~\ref{sec:results} reports a priori and a
posteriori results, cost, and failure modes; Section~\ref{sec:conclusion}
concludes.

% ===================================================================== %
\section{Governing equations}
\label{sec:equations}

\subsection{The QG vorticity equation}
We consider the two-dimensional quasi-geostrophic vorticity equation on a
doubly-periodic domain,
\begin{equation}
\partial_t\omega = \Lop\,\omega + N(\omega),\qquad
N(\omega)=-\Jop(\psi,\omega)+F,\quad \psi=\Delta^{-1}\omega,\quad \dot F=0,
\label{eq:qg}
\end{equation}
with $\Lop=\nu\Delta-\mu-\beta\partial_x\Delta^{-1}$ collecting viscosity,
linear drag, and the $\beta$-plane term, Fourier symbol
\begin{equation}
\hat\Lop(\mathbf k) = -\nu\kappa^2-\mu+i\beta k_x/\kappa^2,\qquad
\kappa=|\mathbf k|,
\end{equation}
and $\Jop(a,b)=\partial_xa\,\partial_yb-\partial_ya\,\partial_xb$. Spatial
discretization is pseudo-spectral with $2/3$ dealiasing of the quadratic
products \citep{canuto2006,boyd2001}; the physical setting follows
\citet{vallis2017}. For the one-step analysis it is convenient to write
\eqref{eq:qg} as an autonomous ODE,
\begin{equation}
\dot\omega = G(\omega) = \Lop\omega + B(\omega,\omega) + F,\qquad
B(a,b)=-\tfrac12\big[\Jop(\Delta^{-1}a,b)+\Jop(\Delta^{-1}b,a)\big],
\label{eq:ode}
\end{equation}
with $B$ symmetric bilinear and $B(\omega,\omega)=-\Jop(\Delta^{-1}\omega,\omega)$,
so $N=B(\omega,\omega)+F$.

\subsection{Exact derivative recursions}
\label{sec:recursions}
Repeated differentiation of \eqref{eq:qg} along the trajectory gives the
mutually recursive relations used throughout:
\begin{equation}
\omega^{(k)} = \Lop\,\omega^{(k-1)} + N^{(k-1)},\qquad
\psi^{(k)}=\Delta^{-1}\omega^{(k)},
\label{eq:wk}
\end{equation}
\begin{equation}
N^{(m)} = -\sum_{j=0}^{m}\binom{m}{j}\,
\Jop\!\big(\psi^{(m-j)},\,\omega^{(j)}\big),\qquad m\ge1,
\label{eq:leibniz}
\end{equation}
where the static forcing drops out of every $N^{(m)}$, $m\ge1$ (bilinearity of
$\Jop$ and the Leibniz rule). Unrolled,
\begin{equation}
\omega^{(k)} = \Lop^k\omega+\sum_{j=0}^{k-1}\Lop^{k-1-j}N^{(j)},\qquad
\text{e.g.}\quad
\Ndot=-\Jop(\dot\psi,\omega)-\Jop(\psi,\dot\omega).
\label{eq:wk-unrolled}
\end{equation}
Each additional time derivative of $N$ costs another layer of nested
Jacobians; this nested-Jacobian cost is what the closure will learn to avoid.
Finally, since $B$ is bilinear, the Fr\'echet expansion of $G$ terminates
exactly:
\begin{equation}
G(\omega+p)=f+Jp+B(p,p),\qquad
f\equiv G(\omega),\quad
J\equiv G'(\omega)=\Lop+2B(\omega,\cdot\,),\quad
G'''\equiv0.
\label{eq:frechet}
\end{equation}

% ===================================================================== %
\section{Time-integration schemes}
\label{sec:schemes}

\subsection{AB2--CN2}
\label{sec:ab2cn2}
Crank--Nicolson treats the stiff linear term implicitly and
Adams--Bashforth-2 the nonlinearity explicitly \citep{ascher1995,frank1997}:
\begin{equation}
\omega^{n+1}=r(h)\,\omega^n+s(h)\Big(\tfrac32 N^n-\tfrac12 N^{n-1}\Big),\qquad
r=\frac{1+\tfrac h2\Lop}{1-\tfrac h2\Lop},\quad s=\frac{h}{1-\tfrac h2\Lop}.
\label{eq:ab2cn2}
\end{equation}
All linear operators are diagonal in Fourier space, so the implicit solve is
an elementwise division; the only nonlinearity is the Jacobian, evaluated
pseudo-spectrally at five FFTs per call with $2/3$ dealiasing. The scheme is
second-order accurate in time.

\subsection{AB4--CN2}
\label{sec:ab4cn2}
The four-level variant replaces the AB2 bracket:
\begin{equation}
\omega^{n+1}=r(h)\,\omega^n+s(h)\,
\frac{55N^n-59N^{n-1}+37N^{n-2}-9N^{n-3}}{24},
\label{eq:ab4cn2}
\end{equation}
with the same $r,s$. Because the CN pair limits the overall order, AB4--CN2 is
still second order; what changes is the \emph{composition} of the truncation
operators (Sec.~\ref{sec:error-imex}) and the explicit stability region
(Sec.~\ref{sec:stability}).

\subsection{RK4}
\label{sec:rk4}
The classical fourth-order Runge--Kutta scheme, fully explicit on $G$:
\begin{equation}
k_1=G(\omega^n),\ \
k_2=G(\omega^n+\tfrac h2k_1),\ \
k_3=G(\omega^n+\tfrac h2k_2),\ \
k_4=G(\omega^n+hk_3),\qquad
\omega^{n+1}=\omega^n+\tfrac h6(k_1+2k_2+2k_3+k_4).
\label{eq:rk4}
\end{equation}
RK4 at a fine step $\hf$ serves two roles: the training-data generator (its
trajectory is the exact flow to $\sim10^{-23}$ at $\hf=10^{-3}$ and below) and
one of the two possible closure targets.

% ===================================================================== %
\section{Analysis of the schemes}
\label{sec:scheme-analysis}

% ------------------------------------------------------------------- %
\subsection{Error analysis}
\label{sec:error}

\subsubsection{Exact one-step map}
Expanding the exact flow in total derivatives (equivalently, the Duhamel
integral in $\varphi$-functions) and substituting \eqref{eq:wk-unrolled},
\begin{align}
\omega(t_n+h)
&= \omega + h(\Lop\omega+N)
 +\tfrac{h^2}{2}(\Lop^2\omega+\Lop N+\Ndot)
 +\tfrac{h^3}{6}(\Lop^3\omega+\Lop^2N+\Lop\Ndot+\Nddot)\notag\\
&\quad+\tfrac{h^4}{24}(\Lop^4\omega+\Lop^3N+\Lop^2\Ndot+\Lop\Nddot+\Ndddot)
\notag\\
&\quad+\tfrac{h^5}{120}(\Lop^5\omega+\Lop^4N+\Lop^3\Ndot+\Lop^2\Nddot
 +\Lop\Ndddot+N^{(4)})+\bO(h^6).
\label{eq:exact-expanded}
\end{align}

\subsubsection{RK4}
Because the Fr\'echet expansion \eqref{eq:frechet} terminates, every RK4 stage
can be expanded \emph{exactly} in the elementary differentials
$\{J^kf,\ B(\cdot,\cdot)\}$; matching against the exact derivatives
\begin{equation}
\omega^{(3)}=J^2f+2B(f,f),\quad
\omega^{(4)}=J^3f+2JB(f,f)+6B(f,Jf),\quad\dots
\end{equation}
shows agreement through $h^4$ and yields the first defect at $h^5$
\citep{hairer1993,butcher2016}.

\begin{proposition}[LTE of RK4]
\label{prop:lte-rk4}
For \eqref{eq:ode},
\begin{equation}
\tau_{\mathrm{RK4}} = h^5\Big[
\tfrac{1}{120}J^4f
-\tfrac{1}{240}J^2B(f,f)
+\tfrac{1}{120}JB(f,Jf)
-\tfrac{1}{60}B(f,J^2f)
+\tfrac{1}{120}B(f,B(f,f))
-\tfrac{1}{80}B(Jf,Jf)\Big]+\bO(h^6)
\equiv h^5T_5+\bO(h^6).
\label{eq:rk4-lte}
\end{equation}
\end{proposition}
\noindent The $J^4f$ coefficient is $1/120$: no stage composition builds the
depth-four chain, so RK4 misses the deepest elementary differential entirely.
The full stage-by-stage derivation is in Appendix~\ref{app:stages}.

\subsubsection{The IMEX pair}
\label{sec:error-imex}
Taylor-expanding the AB brackets
($\tfrac32N^n-\tfrac12N^{n-1}
=N+\tfrac h2\Ndot-\tfrac{h^2}{4}\Nddot+\tfrac{h^3}{12}\Ndddot-\dots$),
the geometric expansions of $r,s$, and subtracting from
\eqref{eq:exact-expanded}:

\begin{proposition}[LTE of AB2--CN2]
\label{prop:lte-ab2}
\begin{equation}
\tau_{\mathrm{AB2}}
=-\frac{h^3}{12}R_3^{(2)}-\frac{h^4}{24}R_4^{(2)}
 -\frac{h^5}{240}R_5^{(2)}-\frac{h^6}{1440}R_6^{(2)}+\bO(h^7),
\label{eq:ab2-lte}
\end{equation}
\begin{align}
R_3^{(2)} &= \Lop^3\omega+\Lop^2N+\Lop\Ndot-5\Nddot,\label{eq:R3}\\
R_4^{(2)} &= 2\Lop^4\omega+2\Lop^3N+2\Lop^2\Ndot-4\Lop\Nddot+\Ndddot,\\
R_5^{(2)} &= 13\Lop^5\omega+13\Lop^4N+13\Lop^3\Ndot-17\Lop^2\Nddot
             +8\Lop\Ndddot-7N^{(4)},\\
R_6^{(2)} &= 43\Lop^6\omega+43\Lop^5N+43\Lop^4\Ndot-47\Lop^3\Nddot
             +28\Lop^2\Ndddot-17\Lop N^{(4)}+4N^{(5)}.
\end{align}
The orders $h^0,h^1,h^2$ cancel identically (second-order accuracy).
\end{proposition}

\begin{proposition}[LTE of AB4--CN2]
\label{prop:lte-ab4}
\begin{equation}
\tau_{\mathrm{AB4}}
=-\frac{h^3}{12}R_3^{(4)}-\frac{h^4}{24}R_4^{(4)}
 -\frac{h^5}{720}R_5^{(4)}+\bO(h^6),
\end{equation}
\begin{equation}
R_3^{(4)}=\Lop^3\omega+\Lop^2N+\Lop\Ndot,\qquad
R_4^{(4)}=2\Lop^4\omega+2\Lop^3N+2\Lop^2\Ndot+\Lop\Nddot,
\end{equation}
\begin{equation}
R_5^{(4)}=39\Lop^5\omega+39\Lop^4N+39\Lop^3\Ndot+24\Lop^2\Nddot
+9\Lop\Ndddot-251\,N^{(4)}.
\end{equation}
\end{proposition}

\begin{remark}[Structural blind spots]
\label{rem:blindspot}
$R_3^{(2)}-R_3^{(4)}=-5\Nddot$: the AB2 stencil is structurally blind to
$\Nddot$, and the weight $-5$ makes $\Nddot$ the dominant, irreducible part of
the AB2--CN2 closure (quantified in Sec.~\ref{sec:error-magnitudes}). AB4--CN2
zeroes the $\Nddot$ slot at orders 3--4 but pays at $h^5$ with the enormous
$-251\,N^{(4)}$ coefficient---the deeper explicit stencil defers, rather than
removes, the nonlinear-derivative burden. Both schemes remain second order:
the leading $\Lop$-monomials are set by the CN pair.
\end{remark}

\begin{remark}[Pad\'e consistency check]
The $\Lop^p\omega$ coefficients of AB2--CN2,
$\{-\tfrac1{12},-\tfrac1{12},-\tfrac{13}{240},-\tfrac{43}{1440}\}$, are exactly
the error coefficients of the $(1,1)$-Pad\'e approximant $r(h)$ to
$e^{h\Lop}$ (pure Crank--Nicolson truncation). All coefficients in
Props.~\ref{prop:lte-ab2}--\ref{prop:lte-ab4} were verified symbolically.
\end{remark}

\begin{remark}[$K$-fine reference]
\label{rem:kfine}
Targeting $K$ fine steps of the same scheme at $\hf=\DT/K$ rather than the
exact flow multiplies the order-$p$ term by $(1-1/K^{p-2})\cdot$(scheme
factors); for the leading term the factor is $(1-1/K^2)$, and for
$K\gtrsim100$ all such factors are $\approx1$. Targeting fine RK4 subtracts
$\tau_{\mathrm{RK4}}$, which first differs from the exact flow at
$\bO(\hf^5)=\bO(\DT^5/K^5)$---negligible. The exact-flow operators are
therefore what one implements in all cases; the choice of target matters only
through the stability margin (Sec.~\ref{sec:closure}).
\end{remark}

\subsubsection{Numerical verification}
\label{sec:lte-verification}
For a developed-flow snapshot $\omega^n$ we form the exact history and future
by fine RK4, take one AB2--CN2 step of size $h$, and measure the empirical
defect and the residuals after subtracting the analytical prediction kept
through order $P$. Keeping $R_3,\dots,R_P$ leaves a residual $\bO(h^{P+1})$;
the measured log--log slopes are $4,5,6$ for $P=3,4,5$, confirming
$R_3,R_4,R_5$ ($R_6$'s slope-7 line sits at the float64 roundoff floor and is
confirmed symbolically instead). Chain-rule derivatives
\eqref{eq:leibniz} are validated against RK4 finite differences to $\le0.7\%$
($\le0.6\%$ for $N^{(5)}$).
\figph{4.2cm}{Figure placeholder: vs-exact-flow convergence.
$\|\tau_{\rm emp}\|$ (slope 3) and residuals after subtracting $R_3$,
$R_3{+}R_4$, $R_3{+}R_4{+}R_5$ (slopes 4,5,6) vs $h$. Source:
\texttt{analysis/validate\_ab2cn2\_vs\_truth.py}.}

% ------------------------------------------------------------------- %
\subsection{Magnitudes, the convergence radius, and error propagation}
\label{sec:error-magnitudes}
The truncation operators are only useful if their magnitudes are known on real
turbulence. Three questions decide the design: how fast do the
$\norm{N^{(m)}}$ grow (convergence radius); how deep may the time stencil
reach (inner wall); and how do estimation errors propagate through the
assembled closure (amplification). All measurements below: float64,
$2/3$-dealiased Jacobians matching the solver right-hand side, $512^2$
forced-turbulence members at marks $\hf=5\times10^{-3}$; truth is the analytic
spectral chain rule \eqref{eq:leibniz}; errors are relative $L_2$ in the
resolved band. Members: FRC-Re25k ($\nu=4.0\times10^{-5}$, $\beta=1$),
FRC-combo ($\nu=4.0\times10^{-5}$, $\beta=0.5$), FRC-kf4
($\nu=1.02\times10^{-4}$, $\beta=1$).

\subsubsection{Outer wall: finite convergence radius}
The defect of one AB2--CN2 step is a power series
$\delta(\DT)=\sum_{p\ge3}c_p\DT^p$, $c_p=R_p/D_p$,
$D_p=\{12,24,240,1440\}$. By Cauchy--Hadamard,
$1/\DTs=\limsup_p\norm{c_p}^{1/p}$; the successive-ratio test only brackets
$\DTs$ (the $c_p$ are not geometric), so the headline number is the root-test
median over $p=3,\dots,6$ (32 anchors per member).

\begin{center}\small
\begin{tabular}{lcclcc}
\toprule
member & $\nu$ & $\beta$ &
$\norm{N^{(m+1)}}/\norm{N^{(m)}}$, $m=0,\dots,4$ &
$\DTs$ (root) & ratio bracket\\
\midrule
Re25k & $4.0\times10^{-5}$ & $1.0$ & $34,\ 58,\ 72,\ 84,\ 93$ & $\mathbf{0.066}$ & $[0.017,0.136]$\\
combo & $4.0\times10^{-5}$ & $0.5$ & $15,\ 27,\ 37,\ 45,\ 51$ & $\mathbf{0.139}$ & $[0.033,0.270]$\\
kf4   & $1.0\times10^{-4}$ & $1.0$ & $10,\ 19,\ 28,\ 33,\ 37$ & $\mathbf{0.199}$ & $[0.043,0.366]$\\
\bottomrule
\end{tabular}
\end{center}

\begin{proposition}[Finite radius]
\label{prop:radius}
$\DTs$ is finite and set by the growth of the $N$-derivatives. For
$\DT>\DTs$ no finite-order analytic closure converges: past $\DTs$, adding
correction orders makes the error \emph{worse}, regardless of how many orders
are carried.
\end{proposition}
\noindent This overturns the naive expectation that higher-order corrections
extend the usable step. Note also that the per-order growth ratios
\emph{increase} with $m$ ($34\to93$ for Re25k): higher derivatives grow faster
than any single rate---derived next.

\subsubsection{The Taylor microscale of the nonlinear term}
\label{sec:microscale}
View $N(t)$ as a stationary process with the domain inner product and define
$\rho(\tau)=\E\langle N(t),N(t+\tau)\rangle/\E\norm{N}^2$. Under stationarity,
$\E\langle N,\Ndot\rangle=\tfrac12\tfrac{d}{dt}\E\norm{N}^2=0$, and
differentiating once more gives
$\E\langle N,\Nddot\rangle=-\E\norm{\Ndot}^2$. Hence the exact identity (no
Gaussianity, no spectral assumption):
\begin{equation}
\boxed{\;
\rho(\tau)=1-\frac{\tau^2}{2}\sigma^2+\bO(\tau^4),\qquad
\sigma^2=\frac{\E\norm{\Ndot}^2}{\E\norm{N}^2},\qquad
\taul=\frac1\sigma.
\;}
\label{eq:microscale}
\end{equation}
$\taul$ is the Taylor microscale of the $N$-process; the same identity holds
for $\omega$ and per wavenumber shell, with
$\sigma(\kappa)=\norm{\dot\omega_\kappa}/\norm{\omega_\kappa}$. Per-order
growth ladders as the moment ratio
\begin{equation}
\sigma_m^2=\frac{\int\sigma(\kappa)^{2(m+1)}E_N(\kappa)\,d\kappa}
                {\int\sigma(\kappa)^{2m}E_N(\kappa)\,d\kappa},
\label{eq:ladder}
\end{equation}
which increases with $m$ because $\sigma(\kappa)$ increases with $\kappa$:
higher derivatives progressively sample the fast small-scale tail---exactly
the escalation ($34\to93$) in the table above. Verification: the identity
holds to $15\%$ on $N$ and $3\%$ on $\omega$ (curvature fit vs $1/\sigma$;
Re25k $\taul\sigma_N=1.15$, kf4 $1.09$; on $\omega$, $0.98$ and $0.97$). Two
further measured facts: the state decorrelates an order of magnitude more
slowly than its nonlinear term ($\sigma_N/\sigma_\omega\approx7.5$--$11$;
over 27 lags $\rho_\omega\ge0.96$ while $\rho_N$ falls to $0.2$--$0.6$), so a
time stencil on the $(\omega,\psi)$-history sees smooth data; and $\rho_N$
plateaus at a positive value ($\approx0.2$, the static forcing plus time-mean
of $-\Jop$), which biases integral decorrelation times but not $\taul$.
\figph{4.0cm}{Figure placeholder: measured $\rho_N(\tau)$, $\rho_\omega(\tau)$
vs the identity parabolas $1-\tau^2\sigma^2/2$ (not fitted), Re25k and kf4.
Source: \texttt{Theoretical\_guarantees/check\_decorrelation\_time.py}
(\texttt{img/rho\_vs\_identity\_*.png}).}

\begin{proposition}[Universal radius law]
\label{prop:two-micro}
$\DTs=C\,\taul$ with $C=2.08\pm0.15$, constant to $7\%$ across $2.5\times$ in
$\mathrm{Re}$ and $2\times$ in $\beta$ (Re25k $2.27$, combo $2.05$, kf4
$1.92$): the analytic closure's convergence horizon is two Taylor microscales
of the nonlinear term. All $(\mathrm{Re},\beta)$-dependence of $\DTs$ enters
through the single measurable rate $\sigma$.
\end{proposition}

\subsubsection{Inner wall: the finite-lag stencil}
\label{sec:inner-wall}
The closure estimates $N$-derivatives from $S$ stored marks at lags
$\{0,\DT,\dots,(S-1)\DT\}$. Depth raises the approximation order
($\bO(\DT^{S-m})$ for the $m$-th derivative) but reaches further back in time;
old marks are stale once the span approaches $\DTs$. Measured ($\Nddot$,
depth 4 vs depth 7, exact spectral spatial operators): at
$\DT=5\times10^{-3}$ depth 7 helps $7.6$--$25\times$ in every member; at
$\DT=1.5\times10^{-2}$ it still helps $3$--$6\times$ for combo/kf4 but
\emph{hurts} for Re25k ($0.99\to1.51$)---precisely where its span
$6\DT=0.09$ exceeds Re25k's $\DTs=0.066$.

\begin{proposition}[Inner wall]
\label{prop:inner}
Deeper stencils lower the $N$-derivative error while the stencil span stays
well inside $\DTs$, and raise it once the span exceeds $\DTs$; the usable step
of a depth-$S$ closure scales as $\DT_{\rm inner}\sim\DTs/(S-1)$. Measured at
Re25k, $S=7$: predicted $0.011$, observed $0.013$.
\end{proposition}
\noindent Consequence for pooled training: a $(\mathrm{Re},\beta)$-blind pool
caps at the \emph{worst} member's inner wall; at $S=7$,
$\DT=1.5\times10^{-2}$ the Re25k samples are unlearnable by any network (the
field decorrelates inside the stencil) and are retained in the pool by design
as past-wall stressors.

\subsubsection{Benign error propagation}
\label{sec:benign}
An error $\varepsilon$ on the learned field enters the closure as
$(\DT^p/D_p)|a_{p,k}|\,\varepsilon\,\norm{\Lop^{p-k}\,\mathrm{field}_k}$;
terms carrying powers of $\Lop$ are potentially amplified. Measured at the
operating step $\DT=5\times10^{-3}$ (per-order errors
$2.6/3.0/4.0\%$ on $\Ndot/\Nddot/\Ndddot$): the un-amplified $\Lop^0\Nddot$
term equals $\norm\delta$ to three figures in every member, and the total
error is $3.00\%$---exactly the $\Nddot$ error. At the envelope edge
$\DT=1.5\times10^{-2}$ (unit per-term errors, pure amplification geometry,
Re25k): $\Lop^0\Nddot$ carries $100.2\%$ of $\norm\delta$, $\Lop^0\Ndddot$
$2.6\%$, $\Lop^1\Ndot$ $0.5\%$, $\Lop^2\Ndot$ $0.00\%$.

\begin{proposition}[Benign amplification]
\label{prop:benign}
Across the entire trained envelope $\DT\le1.5\times10^{-2}$ the closure error
equals the relative $\Nddot$ error, one-to-one and un-amplified. Any
improvement in $\Nddot$ accuracy converts to closure-error reduction 1:1, and
$\Nddot$ accuracy is the single quantity that sets the rollout ceiling.
\end{proposition}
\noindent Two design consequences: the equal-weight per-order relative-$L_2$
loss is the right objective across the envelope (no amplification weighting
needed), and $\Nddot$ validation error is the number to watch.

\begin{corollary}[Operating envelope]
For a depth-$S$ stencil, operate with $(S-1)\DT$ comfortably below
$\DTs(\mathrm{Re},\beta)=2\taul$. Then the series converges
(Prop.~\ref{prop:radius}), the stencil resolves the derivatives
(Prop.~\ref{prop:inner}), and derivative accuracy converts 1:1
(Prop.~\ref{prop:benign}). At $S=7$, Re25k caps the pooled envelope near
$\DT\approx10^{-2}$; combo and kf4 remain inside through
$1.5\times10^{-2}$.
\end{corollary}

% ------------------------------------------------------------------- %
\subsection{Cost analysis}
\label{sec:cost}
Per-step cost is measured in Jacobian evaluations (five FFTs each: two
inverse transforms of $\nabla\psi$, two of $\nabla\omega$, one forward of the
dealiased product; all other operations are Fourier-diagonal multiplies).
\begin{itemize}[leftmargin=1.4em,nosep,topsep=2pt]
\item \textbf{AB2--CN2 / AB4--CN2:} one \emph{new} Jacobian per step; the AB
history reuses cached $N^{n-j}$. Implicit solve: elementwise division. Cost
$\approx5$ FFTs/step.
\item \textbf{RK4:} four Jacobian evaluations per step ($\approx20$
FFTs/step), and no reuse across steps.
\item \textbf{Closed AB2--CN2 (Sec.~\ref{sec:closure}):} the bare step, the
analytical $R_3$ scaffolding (Fourier-diagonal multiplies on the already
available $\hat\omega$, $\hat N$; the $\Lop^3\omega$ piece folds into the
implicit solve at zero cost), and one network forward. Measured in the
production stepper: bare 5 FFTs/step $\to$ closed 8 FFTs/step, plus the
network ($\approx170$ multiply--accumulates per grid point---negligible next
to an FFT). The conditioned variant adds zero further FFTs at rollout: its
spectral context reuses the stepper's own spectral states.
\end{itemize}
To integrate to horizon $T$ at matched accuracy, the bare scheme needs
$T/\hf$ steps at the fine step and the closed scheme $T/\DT$ steps at the
coarse step: the ideal speedup is
\begin{equation}
\text{speedup}\;\approx\;K\times\frac{\text{bare FFTs/step}}
{\text{closed FFTs/step}}
\;=\;K\times\tfrac58
\qquad(\text{vs fine AB2--CN2}),\qquad
K\times\tfrac{20}{8}\quad(\text{vs fine RK4}),
\label{eq:speedup}
\end{equation}
with $K=\DT/\hf$ set at training time. The honest accounting of what fraction
of \eqref{eq:speedup} is realized---given the measured accuracy at each
$\DT$---is in Sec.~\ref{sec:results-cost}.

% ------------------------------------------------------------------- %
\subsection{Stability analysis}
\label{sec:stability}

\subsubsection{Explicit regions}
For the explicit families the absolute-stability region is
$\{z:\ |g(z)|\le1\}$ with $g$ the max-modulus root of
$\rho(\zeta)-z\sigma(\zeta)$ (AB) or $g=R(z)=1+z+\tfrac{z^2}2+\tfrac{z^3}6
+\tfrac{z^4}{24}$ (RK4). For advection the eigenvalues are near-imaginary, so
the binding number is the imaginary-axis stability boundary (ISB): AB2's
imaginary-axis interval is \emph{empty} (marginal at the origin only), AB3's
is $\approx0.724$, AB4's $\approx0.43$, and RK4's is $2\sqrt2\approx2.83$
\citep{hairer1993,butcher2016}. RK4's region is vastly larger per step---but
costs four Jacobians (Sec.~\ref{sec:cost}); AB2 survives in the IMEX pair
because the CN-treated viscosity pulls the advective eigenvalues off the axis.
\figph{3.6cm}{Figure placeholder: absolute-stability regions of AB1--AB4 and
RK4 with axis intercepts. Source: \texttt{analysis/ab\_stability\_regions.py}.}

\subsubsection{IMEX amplification factors}
\label{sec:imex-amp}
On a Fourier mode, split $z=\zL+\zN$ with $\zL=h\hat\Lop(\mathbf k)$ treated
by CN and $\zN$ the explicit advection symbol; the frozen-coefficient
(von Neumann) worst case over background velocity is $\zN=ih\,U|\mathbf k|$.
With
\begin{equation}
r=\frac{1+\zL/2}{1-\zL/2},\qquad \rho=\frac{\zN}{1-\zL/2},
\end{equation}
AB2--CN2 is the two-step recurrence
$\omega^{n+1}=(r+\tfrac32\rho)\,\omega^n-\tfrac12\rho\,\omega^{n-1}$ with
companion amplification
\begin{equation}
g^2-\big(r+\tfrac32\rho\big)g+\tfrac12\rho=0,\qquad
g_{\mathrm{AB2}}=\text{max-modulus root},
\label{eq:g-ab2}
\end{equation}
AB4--CN2 the analogous quartic companion, and
$g_{\mathrm{RK4}}=R(\zL+\zN)$. Absolute stability of the discrete mode is
$|g|\le1$; the maximum stable $\DT$ per regime follows by bisection over all
dealiased modes, using each regime's measured characteristic velocity.
\figph{4.0cm}{Figure placeholder: per-regime stability summary---$|g|=1$
contours at the median $\zL$, the $\zL(\mathbf k)$ cloud, worst-mode $|g|$ vs
$\DT$, and max stable $\DT$ per (regime, scheme) over the 14 ensemble regimes.
Source: \texttt{Theoretical\_guarantees/imex\_stability\_regions.py},
\texttt{analysis/vn\_stability.py}.}

\subsubsection{Truncation symbols}
\label{sec:symbols}
For the closure analysis we need the LTE in amplification space. Replacing
$\Lop^j\omega\to\lambda_L^j$ and $N^{(m)}\to\lambda_N(\lambda_L+\lambda_N)^m$
per mode, the order-$p$ operator becomes the polynomial
\begin{equation}
P_p(\zL,\zN)=a_{p,0}\zL^p+\sum_{k=1}^p a_{p,k}\,\zL^{p-k}\zN(\zL+\zN)^{k-1},
\label{eq:Pp}
\end{equation}
with $a^{(2)}_p$ the AB2--CN2 coefficient rows of
Prop.~\ref{prop:lte-ab2}, so that
\begin{equation}
\hat\tau_{\mathrm{AB2}}=e^z-g_{\mathrm{AB2}}=-\sum_p\frac{P_p}{D_p},\qquad
\hat\tau_{\mathrm{RK4}}=e^z-R(z)=\frac{z^5}{120}+\bO(z^6).
\label{eq:tau-symbols}
\end{equation}

% ===================================================================== %
\section{The step-size resolution closure}
\label{sec:closure}

\subsection{Definition and targets}
Since $\omega^{n+1}_{\rm AB2CN2}+\tau=\omega(t_n+h)$ exactly, the correction
to add to each bare step is $\delta=\tau$, with two natural targets:
\begin{align}
\delta_1 &= \tau_{\mathrm{AB2}}
=-\tfrac{h^3}{12}R_3^{(2)}-\tfrac{h^4}{24}R_4^{(2)}-\tfrac{h^5}{240}R_5^{(2)}
 -\dots
&&\text{(target = exact flow / $K\!\to\!\infty$)},\\
\delta_2 &= \tau_{\mathrm{AB2}}-\tau_{\mathrm{RK4}}
=\delta_1-h^5T_5
&&\text{(target = fine-step RK4)},
\end{align}
with symbols $\hat\delta_1=-\sum_pP_p/D_p$ and
$\hat\delta_2=\hat\delta_1-z^5/120$. Removing the truncation through
$\bO(h^6)$ leaves $\bO(h^7)$:
\begin{equation}
\omega^{n+1}_{\rm closed}=\omega^{n+1}_{\rm AB2CN2}
-\frac{\DT^3}{12}R_3-\frac{\DT^4}{24}R_4-\frac{\DT^5}{240}R_5
-\frac{\DT^6}{1440}R_6 .
\label{eq:closed}
\end{equation}

\subsection{Why target RK4: the stability margin}
\label{sec:margin}
Closing exactly gives $g=g_{\mathrm{AB2}}+\hat\delta_1=e^z+\bO(z^6)$ or
$g=g_{\mathrm{AB2}}+\hat\delta_2=R(z)+\bO(z^6)$. On the advection axis
$z=iy$,
\begin{equation}
|e^{iy}|\equiv1\quad(\text{margin }0),\qquad
|R(iy)|=1-\frac{y^6}{144}+\bO(y^8)<1\quad(\text{margin }>0):
\label{eq:margins}
\end{equation}
RK4 agrees with $e^z$ through $z^4$ and the $z^5/120$ difference is purely
imaginary on the axis (a phase error), so the modulus deficit first appears at
$y^6$---RK4 is faintly dissipative. The exact flow, being neutral, buys no
buffer: any perturbation of an exact-flow-matching closure (analytic tail or
network error) can be destabilizing. This asymmetry, made precise in
Sec.~\ref{sec:nn-stability}, is the deployment rationale for
$\delta_2$/fine-RK4 as the reference: the goal of the whole construction is
AB2--CN2 cost at coarse $\DT$ with the accuracy \emph{and stability} of
fine-$\hf$ RK4.

\subsection{Analytical/learned split}
\label{sec:split}
Each operator separates into a part that is free at inference and a part that
would require nested Jacobians:
\begin{equation}
R_p=\underbrace{R_p^{\rm anal}}_{\text{powers of }\Lop\text{ on }\omega,\,N}
   +\underbrace{R_p^{\rm NN}}_{\text{time derivatives of }N},
\qquad\text{e.g.}\quad
R_3^{\rm anal}=\Lop^3\omega+\Lop^2N,\quad
R_3^{\rm NN}=\Lop\Ndot-5\Nddot.
\end{equation}
The analytical parts are Fourier-diagonal (or reuse the Jacobian already
computed for $N$). The learned parts contain
$\Ndot,\Nddot,\Ndddot,\dots$---each another layer of nested Jacobians
\eqref{eq:leibniz}---and are the only expensive pieces. The network therefore
predicts the \emph{local fields} $[\Ndot,\Nddot,\Ndddot]$; every
$\Lop^k$-weighting (including the nonlocal $\beta$ term) is applied
analytically at assembly, never learned. By Prop.~\ref{prop:benign} the
$\Nddot$ channel dominates the assembled correction 1:1 across the envelope.

\subsection{Inference decomposition}
At inference the linear $\Lop^3\omega$ piece is folded \emph{into} the
implicit solve (it shifts the CN denominator), the $\Lop^2N$ piece is applied
explicitly on the cached Jacobian, and the network supplies the rest:
\begin{equation}
\omega^{n+1}_{\rm closed}
=\underbrace{\text{IMEX}\big[r,s\,;\ \Lop+\tfrac{\DT^2}{12}\Lop^3\big]}
 _{\text{implicit, folded}}
+\underbrace{\DT^3\big(1-\tfrac1{K^2}\big)\tfrac1{12}\,\Lop^2N^n}
 _{\text{explicit, analytical}}
+\underbrace{\DT^3\big(1-\tfrac1{K^2}\big)\,f_{\rm NN}(\text{history})}
 _{\text{explicit, learned}},
\label{eq:inference}
\end{equation}
with $f_{\rm NN}\approx\tfrac1{12}[\Lop\Ndot-5\Nddot]$ the $\DT$-independent
learning target (plus the $R_4$ channel via $\Ndddot$ where warranted). The
network is applied \emph{online}: it is evaluated on the closure trajectory's
own evolving history and its correction feeds the next step, so errors
compound autoregressively and the network never sees truth at deployment.

% ===================================================================== %
\section{Method: the learned estimator}
\label{sec:method}

\subsection{Network architecture}
\label{sec:arch}
The estimator (\emph{CheapDerivClosureNet}) maps the stored step history to
the local derivative fields,
\begin{equation}
X=\big[\omega_0,\dots,\omega_{-(S-1)},\ \psi_0,\dots,\psi_{-(S-1)}\big]
\ \longmapsto\ [\Ndot,\Nddot,\Ndddot],\qquad S=7,
\end{equation}
in four stages that mirror the operator being estimated:
\begin{enumerate}[leftmargin=1.6em,nosep,topsep=2pt]
\item \textbf{Time finite differences (frozen, exact).} Backward-FD rows
$W_{ki}$ from the Vandermonde system
$\sum_iW_{ki}\tfrac{(-i)^p}{p!}=\delta_{kp}$ produce
$\omega^{(k)}_{\rm FD},\psi^{(k)}_{\rm FD}$, $k=0,\dots,3$, via the
per-sample path $W_{ki}/\DT^k$. Nothing is learned here: rescaling the unit
stencil per sample is exact. Accuracy at $S=7$: $\Ndot$ $\bO(\DT^6)$,
$\Nddot$ $\bO(\DT^5)$, $\Ndddot$ $\bO(\DT^4)$.
\item \textbf{Order clip.} Only time-orders $0,\dots,3$ are emitted; the
unused rows (scaled $1/\DT^{4\dots6}\approx10^{10}$--$10^{13}$) are dropped
\emph{before any parameterized layer}. All $S$ snapshots still feed each
emitted row (the order-3 row of a 7-node stencil is 4th-order accurate).
\item \textbf{Learnable spatial gradients.} Depthwise width-15 stencils
(14th-order central-difference initialization) apply $\partial_x,\partial_y$
to the 8 differentiated fields, circular padding. The stencil parameter is the
\emph{dimensionless} unit-spacing operator; the physical $1/dx,1/dy$ scaling
is applied per sample at forward---exact by linearity, and the optimizer never
sees the $1/dx$ amplification. One stencil pair serves all grids.
\item \textbf{Jacobian features and physics-initialized mix.} The
$(4)^2=16$ pointwise products
$\Jop(\psi^{(i)},\omega^{(j)})$ feed a $1\times1$ mix initialized to the exact
chain-rule binomials of \eqref{eq:leibniz} ($w[m-1,(m-j,j)]=-\binom mj$): at
initialization the network \emph{is} the analytic estimator up to the two FD
truncations.
\end{enumerate}
Totals: $3{,}700$ parameters, of which the 49 time-FD weights are an inert
frozen buffer (the trained capacity is the stencils and the mix);
$\approx170$ multiply--accumulates per
grid point; no FFTs inside the model---the pointwise product is the sole
irreducible nonlinearity, and carrying $\psi$ in the history makes even the
$\beta$ term local ($\beta\partial_x\Delta^{-1}\omega=-\beta v$). The
architecture is exactly quadratic in the input field, matching the quadratic
nonlinearity of \eqref{eq:ode}.

\begin{remark}[Why the spatial stencils are learnable at all]
Their job is double: approximate $\partial_x,\partial_y$ \emph{and} absorb the
differentiated time-FD error. The second job is the Wiener mechanism analyzed
next, and is why the width was set to 15 and why the stencils are conditioned
in Sec.~\ref{sec:conditioning}.
\end{remark}

\subsection{Exact error analysis of the estimator}
\label{sec:wiener}
Two error operators separate cleanly. The time-FD estimate carries
\begin{equation}
\varepsilon_k\equiv\omega^{(k)}_{\rm FD}-\omega^{(k)}
=C_k\,\DT^{S-k}\,\omega^{(S)}+\bO(\DT^{S-k+1}),\qquad
C_k=\frac1{S!}\sum_{i=0}^{S-1}W_{ki}(-i)^S,
\label{eq:eps-k}
\end{equation}
with $\varepsilon_0=0$ and $\varepsilon^\psi_k=\Delta^{-1}\varepsilon_k$
($\Delta^{-1}$ commutes with the stencil). The learned gradients carry
$D_\theta=\partial+\delta_\theta$ with $\delta_\theta$ the stencil gap.
Propagating both through the bilinear Jacobian and discarding
$\bO(\varepsilon\delta,\varepsilon^2,\delta^2)$,
\begin{equation}
\Delta N^{(m)}=T^{(m)}+\mathcal A^{(m)}[\delta_\theta],\qquad
T^{(m)}=\DT^{S-m}C_m\Big[-\Jop\big(\Delta^{-1}\omega^{(S)},\omega\big)
-\Jop\big(\psi,\omega^{(S)}\big)\Big]+\bO(\DT^{S-m+1}):
\label{eq:DeltaN}
\end{equation}
the leading time error \emph{at every order $m$} is driven by the same field
$\omega^{(S)}$, entering through exact Jacobians, while
$\mathcal A^{(m)}[\delta_\theta]$ is linear in the stencil gap. Training
therefore solves a Wiener problem \citep{wiener1949}: the
least-squares-optimal stencil correction $\delta^\star$ pre-cancels the time
error it can reach.

\subsubsection{The pooled floor and its three parts}
For a pool of tiers $j=(\text{member},\DT)$ with weights $w_j$, orthogonal
decomposition of the quadratic loss about the per-tier optima gives
\begin{equation}
\mathcal L^{\rm uncond}_{\min}
=\underbrace{\sum_jw_j\,\mathcal L_j(\delta^{\star\star}_j)}
  _{\text{(i) cascade remainder}}
+\underbrace{\min_{\bar\delta}\sum_jw_j\E_j
  \norm{\mathcal A[\delta^\star_j-\bar\delta]}^2}
  _{\text{(ii) pooled-variance mismatch}}
+\underbrace{\sum_jw_j\E_j
  \norm{\mathcal A[\delta^{\star\star}_j-\delta^\star_j]}^2}
  _{\text{(iii) width-}W\text{ shape deficit}},
\label{eq:floor}
\end{equation}
where $\delta^\star_j$ is the best width-$W$ filter for tier $j$ and
$\delta^{\star\star}_j$ the best shell-diagonal transfer. Measured at
$S=7$, $m=2$, $W=15$ (valid band): diagonal-absorbable fraction
$0.21$--$0.29$; full-slot absorbable $0.36$--$0.44$;
$\text{(ii)}=0.72$; $\text{(iii)}=0.31$. A single unconditioned stencil can
only fit the pooled mean; its net absorption is $\approx0.03$---the pooled
variance (ii) dominates everything. Deleting (ii) is what conditioning is for.

\subsubsection{Single-mode analysis of the driver}
Under a sweeping split $\hat N=-i(\mathbf k\cdot\mathbf U)\hat\omega+\hat f$
(A1), frozen coefficients over the stencil span (A2), advective band
$\Omega^2\gg\lambda^2$ (A3), and subdominant triads in the shell balance (A4),
induction on $\hat\omega^{(n)}=a^n\hat\omega+\sum_ja^j\hat f^{(n-1-j)}$,
$a=i\Omega+\lambda$, gives the per-shell transfer between derivative orders:
\begin{equation}
\widehat{\omega^{(S)}}_\kappa
=\sigma_\omega(\kappa)^{S-k}e^{i\phi(\kappa)}\,\widehat{\omega^{(k)}}_\kappa
+\hat r,\qquad
\phi=(S-k)\tfrac\pi2\,\mathrm{sign}(\Omega),\qquad
\hat r=\hat r_1+\hat r_2+\hat r_3,
\label{eq:transfer}
\end{equation}
with the shell rate
$\sigma_\omega(\kappa)^2=\langle\Omega^2\rangle_{\rm shell}
=\tfrac12\kappa^2U^2+\tfrac12\beta^2/\kappa^2-U\beta\cos\alpha_{\mathbf U}$.
The remainder taxonomy bounds what any conditioning can recover:
$r_1$ (cascade input, shell-incoherent with $\omega^{(k)}$;
\emph{irreducible}---measured incoherent fraction $0.71$--$0.79$ on healthy
tiers); $r_2$ (coefficient drift within the stencil span, size
$\tfrac{(S-k)(S+k-1)}2\,\tau_\kappa/T_U$ with $\tau_\kappa=1/\sigma_\omega(\kappa)$
the per-shell state timescale and $T_U$ the large-eddy turnover, conditionable
in time); $r_3$ (intra-shell
anisotropy: for odd $S-k$ the sweeping transfer $\propto\cos^{S-k}\theta$
shell-averages to \emph{zero}, so an isotropic correction captures none of it
and only $\cos3\theta,\cos5\theta$ harmonics are reachable by angular
features).

\subsection{Physics conditioning}
\label{sec:conditioning}
Combining \eqref{eq:eps-k} and \eqref{eq:transfer}, the per-shell
regression coefficient of the time error on the resolved field factorizes:
\begin{equation}
r^{(k)}_j(\kappa)
=\underbrace{\DT_j^{\,S-k}}_{\text{exact}}\times
\underbrace{C_k\,\sigma_\omega(\kappa;\mathrm{Re},\beta,\mu)^{S-k}
e^{i\phi(\kappa)}}_{g^{(k)}(\kappa)}+\text{h.o.t.}
\label{eq:factorization}
\end{equation}
Two structural facts follow. First, the optimal stencil correction varies as
$(\DT\,\sigma(\kappa))^{S-k}$---a factor $3^5\approx243$ across our
$\DT$-sweep at $k=2$---which \emph{derives} the physics conditioning: no
single unconditioned stencil can track it. Second,
$(\mathrm{Re},\beta,\mu)$ enter only through $\sigma_\omega(\kappa)$, so the
regime never needs to be supplied as scalars: it is measurable from the
model's own input stack. The two-mark estimator with arcsin de-biasing (a
single mode rotating at $\Omega$ has FD response
$\tfrac2\DT|\sin(\Omega\DT/2)|$),
\begin{equation}
\shat(\kappa)=\frac2\DT\arcsin\!\Big(\min\big(
\tfrac\DT2\,\tfrac{\norm{(\hat\omega_0-\hat\omega_{-1})_\kappa}}
{\DT\,\norm{(\hat\omega_0)_\kappa}},\,1\big)\Big),
\label{eq:sigmahat}
\end{equation}
recovers $\sigma(\kappa)$ with median relative error $0.000$--$0.001$:
data-conditioning is lossless. The measured $\DT$-collapse of
$r_j/\DT^{S-k}$ across the sweep is $1.6\times$ (vs $243\times$ raw),
confirming \eqref{eq:factorization}.

\paragraph{Conditioned architecture (production: local taps).}
The deliverable keeps the zero-FFT locality of the estimator and conditions
the stencils by \emph{tap modulation}:
\begin{equation}
\mathrm{taps}^{(k)}_d
=\mathrm{taps}^{(k)}_{d,\rm base}
+\Big(\frac{\DT}{\DT_{\rm ref}}\Big)^{S-k}
\Delta\mathrm{taps}_{\theta,k,d}\big(\shat\text{-features}\big),
\qquad k\ge1,
\label{eq:condlocal}
\end{equation}
with $\DT_{\rm ref}=1.5\times10^{-2}$ anchoring the exact scaling law at
$\bO(1)$ (raw $\DT^{S-k}\approx10^{-11}$--$10^{-16}$ renders the path
numerically dead: raw $\DT^{S-k}\lesssim5\times10^{-8}$ down to
$\sim10^{-14}$ across the sweep), amplitude exactly zero on the $k=0$ channels
($\varepsilon_0=0$: the undifferenced mark has no time error to absorb), and
$\Delta\mathrm{taps}_\theta$ a small shared MLP
($10\to24\to24\to240$, $\tanh$, $\approx6.9$k parameters) reading
$[x_i,\log(1{+}x_i)]$, $x_i=\DT\,\shat(\kappa_i)$, at five fixed relative
shells $\kappa_i/\kappa_{\max}\in\{0.15,0.3,0.5,0.7,0.85\}$. The output head
is zero-initialized: at initialization the conditioned model is bit-identical
to the unconditioned control. A spectral variant
$\hat D^{(k)}_d=ik_d[1+\DT^{S-k}g^A_\theta(x,\tilde\kappa)]
+k_d\DT^{S-k}g^B_\theta(x,\tilde\kappa)$ has zero shape deficit and serves as
the ceiling instrument, but costs $\approx24$ FFTs per forward and is not
deployed.

\paragraph{Predicted ceilings.}
The conditioned floor deletes term (ii) of \eqref{eq:floor}; per-tier,
$\mathrm{floor}^{\rm cond}_j
=\mathrm{raw}_j\sqrt{\mathrm{irr}_j+\mathrm{coh}_j(\mathrm{def}(\mathcal F)
+\text{h.o.t.})}$. Predictions: kf4@$1.5\times10^{-2}$: $0.031\to0.023$;
FRC-256@$1.5\times10^{-2}$: $0.047\to0.037$; FRC-256@$10^{-2}$:
$0.0068\to0.0055$; pooled raw floor over 9 tiers $0.175$ (vs the observed
unconditioned plateau $0.19$); conditioned-pooled excluding the past-wall tier
$\approx0.05$. Section~\ref{sec:results-apriori} confronts these with the
trained result.

\subsection{Loss}
\label{sec:loss}

\paragraph{Offline derivative loss.}
Per-sample, per-order relative $L_2$ with equal weights over
$[\Ndot,\Nddot,\Ndddot]$ (Prop.~\ref{prop:benign} makes equal weights
correct), predictions $2/3$-dealiased before the loss, and a norm floor
$\max(\norm{t_c},\ 0.1\times\text{member-median}\norm{t_c})$ in the
denominator to cap the leverage of near-zero-target samples without dropping
them. Medians are logged alongside floored means: relative errors on
near-quiescent samples otherwise poison the optimizer (Sec.~\ref{sec:training}).

\paragraph{Rollout fine-tune.}
The offline objective does not see the closed loop, so the estimator is
fine-tuned \emph{through the stepper}: the supervised head is
$\frac1M\sum_{m=1}^M\mathrm{rel}L_2(\omega^{\rm closed}_m,\omega^{\rm truth}_m)$
against fine-RK4 truth over $M$ coarse steps (truncated backpropagation,
window 4; $M$ scheduled up to 21), followed by a truth-free tail. The model
that produced every a posteriori result of Sec.~\ref{sec:results} used, on
that tail, an \emph{annulus-enstrophy hinge}
$\mathrm{relu}\big(\log Z_{\rm ann}^{(f)}/Z_{\rm ann}^{(f-1)}\big)$ over the
aliased band $|\mathbf k|>\tfrac23 k_{\max}$, weight $10^{-3}$ (larger
weights destabilize the supervised head---an objective-imbalance failure we
observed directly).

\paragraph{Upgraded tail (current training campaign; results pending).}
Two replacements for the hinge are in a running $\lambda$-sweep whose results
had not landed at the time of writing and are \emph{not} claimed in
Sec.~\ref{sec:results}:
\begin{enumerate}[leftmargin=1.6em,nosep,topsep=2pt]
\item \emph{Analytic free-tail supervision.} On truth-free rolled steps the
model heads are supervised against the on-the-fly \emph{exact analytic}
targets \eqref{eq:leibniz} of the pre-step rolled state (weight
$10^{-2}$, capped): the estimator must remain a consistent derivative
estimator on its own trajectory, arbitrarily far from the data horizon.
\item \emph{Von Neumann certificate penalty} (derived in
Sec.~\ref{sec:nn-stability}):
$\mathcal L_{\rm stab}
=\lambda\,\mathrm{mean}_\kappa\,\mathrm{relu}\big(|G_{\rm eff}(\kappa)|
-(1-\epsilon)\big)^2$, $\epsilon=0.02$, restricted to the linearization's
validity shells $|\DT\,\shat(\kappa)|\le\tfrac12$;
$\lambda\in\{0.1,1,10\}$ is the only new hyperparameter.
\end{enumerate}

% ===================================================================== %
\section{Stability analysis of the learned closure}
\label{sec:nn-stability}

\subsection{Amplification of the closed scheme}
The correction enters the update additively, hence the amplification
additively:
\begin{equation}
g_{\mathrm{NN},i}(z)=g_{\star,i}(z)+\hat\epsilon_i(z),\qquad
\hat\epsilon_i=\hat\epsilon_{\mathrm{an},i}+\hat\epsilon_{\mathrm{NN},i},
\label{eq:gNN}
\end{equation}
with $g_{\star,1}=e^z$, $g_{\star,2}=R(z)$, $\hat\epsilon_{\rm an}$ the
deterministic leftover from cropping the closure at order $p$, and
$\hat\epsilon_{\rm NN}$ the stochastic fit residual.

\begin{proposition}[Finite-order closure: the base shows through]
\label{prop:base-shows}
Assembling $R_3,\dots,R_p$ cancels the base scheme's truncation series
term-by-term up to order $p$ only:
\begin{equation}
g_{\mathrm{NN},i}
=g_{\star,i}
\underbrace{-\,\tau^{\rm base}_{p+1}z^{p+1}+\bO(z^{p+2})}
_{\hat\epsilon_{\mathrm{an},i}}
+\underbrace{\epsilon_{\rm rel}\,\bO(h^3)}_{\hat\epsilon_{\mathrm{NN},i}}.
\end{equation}
The surviving tail is the \emph{base scheme's own} ($\tau^{\rm AB2}_k\neq
\tau^{\rm AB4}_k$; e.g.\ $R_3$ carries $-5\Nddot$ for AB2 and $0$ for AB4),
never the target's: above the assembled order the closure reverts to
$g_{\rm base}$, even when matching RK4.
\end{proposition}

\begin{corollary}[Empirical closure, $K\to\infty$]
The full-tail empirical correction $\delta=\Phi_{\rm ref}-\Phi_{\rm base}$ is
a difference of maps, not a truncated series: $\hat\epsilon_{\rm an}=0$ and
$g_{\mathrm{NN},i}=g_{\star,i}+\epsilon_{\rm rel}\bO(h^3)$, base-agnostic.
\end{corollary}

\subsection{Margin bound}
By the triangle inequality, $|g_{\mathrm{NN},i}|\le1$ is guaranteed whenever
$\epsilon_i\le m_i(z)=1-|g_{\star,i}(z)|$, tight on the advection axis. There
the two targets differ qualitatively (cf.\ \eqref{eq:margins}):
$m_1(iy)=0$---no perturbation of an exact-flow-matching closure can be
certified stable---while
\begin{equation}
\boxed{\,m_2(iy)=\tfrac{y^6}{144}+\bO(y^8)\,}:
\end{equation}
matching RK4, the closed scheme is provably stable whenever
$\big|{-}\tau^{\rm base}_{p+1}(iy)^{p+1}\big|+\epsilon_{\rm rel}\bO(h^3)
\le y^6/144$. Refining $h$ helps the network term; the analytic-ceiling term
is an $h$-independent comparison of $y$-coefficients set by the base scheme.
This margin asymmetry is the quantitative form of the deployment rationale of
Sec.~\ref{sec:margin}.

\subsection{Accumulated error}
Over $n=T/h$ steps,
$\norm{e(T)}\le\epsilon_i\sum_{m<n}|g|^m$, which for $|g|\le1$ gives
\begin{equation}
\text{match exact: }
\norm{e(T)}\le C_{\rm base}h^pT+\epsilon_{\rm rel}\bO(h^2)T;
\qquad
\text{match RK4: }
\norm{e(T)}\le Ch^4T+C_{\rm base}h^pT+\epsilon_{\rm rel}\bO(h^2)T,
\end{equation}
(drop $C_{\rm base}h^pT$ in the empirical case). The learned floor
$\epsilon_{\rm rel}\bO(h^2)T$ preserves second-order convergence, is linear in
horizon and in network quality; the finite-order floor is smaller in $h$ but
frozen (reducible only by adding analytic orders) and base-dependent.

\subsection{A trainable certificate}
\label{sec:certificate}
The margin bound needs $|g_{\rm NN}|$ per mode \emph{for the actual trained
network}. Linearizing the learned piece about a frozen advecting background
(the Wiener linearization: Jacobians frozen to $i\sigma$ advection, the
stencil pair entering through its trigonometric-polynomial symbol ratio
$\rho(\kappa)=\hat W_{\rm base+mod}(\theta)/\hat W_{\rm exact}(\theta)$
evaluated on the sample's own $\shat$ shells) yields the closed AB2--CN2
companion with effective symbols
\begin{equation}
I=\hat\Lop+\tfrac{\DT^2}{12}\hat\Lop^3\ \ (\text{implicit, folded}),\qquad
E=i\shat\big(1+\tfrac{\DT^2}{12}\hat\Lop^2\big)
+\tfrac1{12}\big(\hat\Lop\,T_1-5\,T_2\big)\ \ (\text{explicit}),
\end{equation}
$T_c$ the linearized transfer of the learned $\Ndot,\Nddot$ heads, and
\begin{equation}
G_{\rm eff}(\kappa)=\text{spectral radius of }
\begin{pmatrix} a+b & c\\ 1 & 0\end{pmatrix},\qquad
a=r\big(1+\tfrac\DT2 I\big),\ \ b=\tfrac32\DT rE,\ \ c=-\tfrac12\DT rE,\ \
r=\frac1{1-\tfrac\DT2 I}.
\label{eq:geff}
\end{equation}
$G_{\rm eff}$ is differentiable in every trainable tensor (taps, modulation,
mix), reproduces the analytic-arm stability exactly at zero modulation
($\rho\equiv1$), and is evaluated only on shells with
$|\DT\shat|\le\tfrac12$, outside which the frozen-coefficient model provably
departs from the dynamics and the penalty would be noise. It serves as both
diagnostic (a per-step contraction certificate is \emph{horizon-free},
removing the fixed-horizon asterisk of rollout training structurally) and
penalty ($\mathcal L_{\rm stab}$ of Sec.~\ref{sec:loss}). Caveats: it is
linearized (the $r_2,r_3$ remainders are neglected) and isotropic per shell.
On the current stable reference model the certificate reads
$|G_{\rm eff}|$ valid-shell mean $0.9992$, max $0.9997$---marginal at large
$\DT$, as it must be.

% ===================================================================== %
\section{Ensemble training}
\label{sec:training}

\paragraph{Data.}
Training data are deep 28-mark windows advanced by fine RK4 from developed
snapshots of an ensemble of forced-turbulence members spanning the
$(\mathrm{Re},\beta)$-range ($512^2$, FRC-256 at $256^2$; domains $4\pi$,
isotropic), sliced into $S=7$ stencils at the sweep steps
$\DT\in\{5\times10^{-3},10^{-2},1.5\times10^{-2}\}$. The unconditioned
control pool is 18 (member,\,$\DT$) roots
(FRC-\{b2, b25, kf4, Re25k, combo, 256\}$\times$3 steps); the conditioned run
trained on an enlarged 41-root pool that adds further forced members
(FRC-\{b0, b05, b075, b1\}) and decaying-turbulence members. The hygiene
control of Sec.~\ref{sec:results-apriori} isolates that data difference on
the 17 shared roots. Targets $[\Ndot,\Nddot,\Ndddot]$ are computed
once per window by the exact spectral chain rule \eqref{eq:leibniz}
(forcing rebuilt per member from its manifest; the targets themselves are
forcing-free since $m\ge1$). The $\DT$ sweep deliberately includes
near-wall and past-wall tiers (Re25k at $1.5\times10^{-2}$ exceeds its
$\DTs$).

\paragraph{Pipeline safeguards (each one earned by a failure).}
\begin{itemize}[leftmargin=1.4em,nosep,topsep=2pt]
\item \emph{By-window splits.} Whole windows go to one split; anchored slices
within a window are near-duplicates $\sim\hf$ apart, so per-sample random
splits leak and chronological splits shift covariates.
\item \emph{Quiescent-window filter.} Early quasi-zonal spin-up flow has
$\Jop(\psi,\omega)\approx0$ and derivative targets $10^4$--$10^5\times$
smaller than developed flow; per-sample relative losses explode on them and
the optimizer's cheapest move is shrinking \emph{all} predictions toward
zero. Windows are dropped when the target norm falls below $10^{-2}\times$
the member median or the stack roughness freezes. Spin-up duration is
$\beta$-dependent (9--33\% of windows across members), so per-member start
times are set by a developed-flow criterion, and the filter is the safety
net.
\item \emph{Multigrid handling.} Batches are shape-homogeneous but the
$dx$-rescale of the stencils is strictly per-sample (equal shape does not
imply equal $dx$); dealias masks are per-shape; anisotropic members are
refused.
\item \emph{Precision.} All closure-pipeline compute is float64: the raw
target is $\bO(\DT^3)\approx10^{-9}$, below float32 machine epsilon. (Disk
storage of inputs is float32, upcast at load; this floors only the deepest
channel $\Ndddot$, which carries 2.6\% of the closure.)
\end{itemize}

\paragraph{Optimization.}
Offline: AdamW, lr $5\times10^{-5}$, weight decay $10^{-4}$, cosine schedule,
200 epochs, batch 4, float64, physics init, seed 0. Rollout fine-tune
(reference run, whose epoch-33 checkpoint produced all a posteriori results):
warm start from the offline conditioned model; strides $1,2,3$
($\DT=\mathrm{stride}\times5\times10^{-3}$); unroll schedule
$4{:}6$, $8{:}8$, $12{:}8$, $16{:}8$, $21{:}10$ with per-stride clamps
$21/7/3$; truncated BPTT window 4; lr $5\times10^{-5}$; grad-clip 1.0;
batch 1 with member-pure accumulation; 40 epochs.

% ===================================================================== %
\section{Results}
\label{sec:results}

\subsection{A priori accuracy and the floors}
\label{sec:results-apriori}
Pooled test relative $L_2$ per order (means; medians in parentheses where
they differ materially):
\begin{center}\small
\begin{tabular}{lcccc}
\toprule
model & pooled & $\Ndot$ & $\Nddot$ & $\Ndddot$\\
\midrule
unconditioned control (\texttt{deriv7\_filtered\_floor0.1})
 & 0.269 & 0.058 & 0.186 & 0.563\\
conditioned (\texttt{cond\_local}, best ep63)
 & 0.214 & 0.074\,(0.054) & \textbf{0.138}\,(0.093) & 0.430\,(0.178)\\
hygiene control (equal 17-root data, ep107)
 & 0.250 & 0.059 & 0.178 & 0.511\\
\bottomrule
\end{tabular}
\end{center}
Three readings. (a) The unconditioned $\Nddot$ plateau $0.186$ matches the
predicted pooled raw floor $0.175$ of \eqref{eq:floor}: the control is
data-efficient but variance-limited, as derived. (b) Conditioning is real on
equal data: the conditioned model beats the hygiene control on $\Nddot$ on
16/17 shared roots (e.g.\ kf4 across the three $\DT$:
$0.058/0.065/0.057$ vs $0.085/0.088/0.086$), and the improvement concentrates
at low $\kappa$ (the $\varepsilon(\kappa)$ profile is U-shaped, not a high-$k$
cliff). (c) \emph{The predicted conditioned ceilings were not reached}: kf4@
$1.5\times10^{-2}$ plateaued at $0.065$ vs the predicted $0.023$, pooled
$0.214$ vs $\approx0.05$; the run was stopped at plateau (epoch 78 of 300).
The prediction is open, not falsified---the gap is the difference between the
information-theoretic floor of the conditioned \emph{family} and what this
optimizer/parameterization extracted---but we report it as unmet.
\figph{4.0cm}{Figure placeholder: per-(member,$\DT$,order) a priori rel-$L_2$
heatmap, conditioned vs control, with the per-tier Wiener floors overlaid.
Source: \texttt{training/eval\_deriv\_by\_root.py},
\texttt{Theoretical\_guarantees/check\_wiener\_floor.py}.}

\subsection{A posteriori rollout: accuracy}
\label{sec:results-apost}
Protocol: 16 coarse closed steps against fine-RK4 truth
($\hf=10^{-5}$; $K=500$--$1500$), improvement $=$
$\mathrm{rel}L_2^{\rm bare}/\mathrm{rel}L_2^{\rm closed}$ at $t=16\DT$; 10
initial-condition draws per $\DT$ across three members---kf4 and combo at
$512^2$, FRC-256 at $256^2$---with combo held out of the fine-tune entirely.

\paragraph{Analytic ceiling.}
The network-free $R_3$ closure (exact chain-rule $\Ndot,\Nddot$; no learned
component) is stable in every arm and sets the honest upper bound:
$71.4\times$ / $35.1\times$ / $132.6\times$ over bare at
$\DT=5\times10^{-3}/10^{-2}/1.5\times10^{-2}$. A one-step
truncation-ratio computation predicts $66.6/30.7/19.1\times$: the two agree
within 7--14\% at the two smaller steps, while the $132.6\times$ at
$1.5\times10^{-2}$ exceeds its one-step ceiling because the \emph{bare}
baseline there enters super-linear error growth---an accumulation property of
the baseline, not of the closure.

\paragraph{Learned closure.}
After rollout fine-tuning, over the 10-draw replication:
\begin{center}\small
\begin{tabular}{lccc}
\toprule
$\DT$ & stable & improvement over bare, median [min, max] & vs analytic ceiling\\
\midrule
$5\times10^{-3}$   & 10/10 & $\mathbf{15.96\times}$ [0.92, 23.57] & $71.4\times$\\
$10^{-2}$          & 9/10  & $0.50\times$ [0.01, 6.17]            & $35.1\times$\\
$1.5\times10^{-2}$ & 6/10  & $0.65\times$ [0.33, 1.89]            & $19.1\times$ (one-step)\\
\bottomrule
\end{tabular}
\end{center}
The $16\times$ gain at the operating step replicates across draws and
members, including the held-out member (its worst draw is parity, $0.92$).
At the two larger steps the learned closure is at or below parity where it is
stable, and stability itself degrades to 60\% of draws at
$1.5\times10^{-2}$. For calibration: the pre-fine-tune model blew up at the
two larger steps (steps 7 and 13) and was stable but \emph{below} parity
($0.72\times$) at the operating step---the fine-tune converted
$0.72\times$ into $16\times$ there and only partially cured the envelope
edge. The gap between $16\times$ and the $71\times$ analytic ceiling
is the direct cost of the remaining $\Nddot$ estimation error
(Prop.~\ref{prop:benign}).
\figph{4.6cm}{Figure placeholder: rollout error growth, bare vs closed vs
analytic-$R_3$, per $\DT$; vorticity fields and differences at $t=16\DT$;
improvement-factor box plots over the 10 draws. Source:
\texttt{diagnostics/rollout\_aposteriori.py},
\texttt{diagnostics/consolidate\_apost\_cases.py}.}

\subsection{A posteriori rollout: stability, the horizon wall, and the mechanism}
\label{sec:results-stability}
Three stability findings, reported in full.

\paragraph{The instability is network-specific.}
The analytic $R_3$ arm is stable at every $(\DT,\mathrm{IC})$ tested, at CFL
up to $0.85$; the learned arm can diverge even at CFL $0.56$. Removing the
$\Nddot$ channel from the learned correction restores stability in every arm
but also removes all gain ($1.00\times$ everywhere): the destabilizer and the
value live in the same channel.

\paragraph{Mechanism: history contamination.}
The blow-up is a feedback loop specific to the memory of the method: network
noise enters the correction, contaminates the stored 7-level history, is
amplified by the $1/\DT^k$ time-FD weights into noisier derivative estimates,
and returns as a noisier correction. The blow-up horizon tracks (steps until
injected noise $\approx$ truncation signal) $+\,{\sim}4$. Aliasing is not the
carrier (the learned arm gets \emph{worse} in a fully dealiased world;
the analytic arm is indifferent), and the radial enstrophy-injection
component is small and $\DT$-independent---a dissipative projection of the
correction (P0 experiment) shifts blow-ups by $\le\pm1$ step and fails its
pre-registered stability bars, ruling out the cheap inference-time fix. The
within-shell (phase) component identified by the certificate is the active
target.

\paragraph{The horizon wall.}
Extending the $5\times10^{-3}$ arm from 16 to 64 and 128 steps: \emph{0/10
draws stable in both arms}, blow-ups at steps 21--38, i.e.\ at
$\approx2\times$ the fine-tune's free horizon. The replicated $16\times$ gain
is therefore horizon-limited: valid within the trained window, not
unconditional. This is the sharpest open problem; the certificate penalty
(Sec.~\ref{sec:certificate}) attacks exactly this---a per-step contraction
certificate is horizon-free by construction---and the corresponding
$\lambda$-sweep is running at the time of writing.
\figph{3.6cm}{Figure placeholder: blow-up step vs horizon, 10 draws, 16/64/128
steps, with the free-horizon line; certificate $|G_{\rm eff}(\kappa)|$
histograms before/after the penalty. Source:
\texttt{diagnostics/aposteriori\_stability.py},
\texttt{training/wiener\_certificate.py}.}

\subsection{Cost}
\label{sec:results-cost}
Measured on the production stepper ($512^2$): the closed step costs 8 FFTs vs
the bare 5, plus $\approx170$ MACs/point of network---a per-step overhead of
$\approx60\%$, against the $K$-fold step-count reduction of
\eqref{eq:speedup}. At the replicated operating point
($\DT=5\times10^{-3}$, $16\times$ error reduction on-horizon), matching the
closed accuracy with the bare scheme would require $\approx4\times$ smaller
steps (second-order scheme), so the realized speedup at matched accuracy is
$\approx(4\times5)/8=2.5\times$ against bare AB2--CN2 on the trained horizon;
the ideal bound \eqref{eq:speedup} is approached only as the learned
$\Nddot$ error approaches the analytic ceiling.
The conditioned context costs zero additional FFTs at rollout (spectral
states reused); training-side costs are modest (rollout fine-tune
$\approx2$--$3$ GPU-h; the offline ensemble run $\approx20$ GPU-h; truth
generation dominates and is cached across all experiments).
\TODO{tighten the realized-speedup accounting with wall-clock numbers from
\texttt{rollout\_timed\_pareto.py} on the final model.}

% ===================================================================== %
\section{Conclusions and future work}
\label{sec:conclusion}
We posed the time-discretization error of spectral AB2--CN2 as a closure
problem and carried it from closed-form truncation theory to a deployed,
certificate-regularized learned scheme. The theory side is complete and
verified: LTEs of the scheme family to $\bO(h^6)$
(Props.~\ref{prop:lte-rk4}--\ref{prop:lte-ab4}), a finite and universal
convergence radius $\DTs=2.08\,\taul$ (Prop.~\ref{prop:two-micro}), an inner
stencil wall (Prop.~\ref{prop:inner}), benign 1:1 error propagation
(Prop.~\ref{prop:benign}), the margin asymmetry that makes fine-RK4 the right
target (Sec.~\ref{sec:margin}), and an exact Wiener analysis of the estimator
that derives---and prices---physics conditioning
(Sec.~\ref{sec:wiener}). The practice side is honest: the analytic $R_3$
closure alone is unconditionally stable in our tests and worth up to two
orders of magnitude; the learned closure replicates $16\times$ at the
operating step on the trained horizon, beats its matched control on the
dominant $\Nddot$ channel on 16 of 17 shared roots as an estimator,
but has not reached its predicted conditioned floor, is at
parity at the envelope edge, and remains horizon-limited by a diagnosed
history-contamination loop.

\paragraph{Future work.}
(i) \emph{Certificate-regularized training} ($\lambda$-sweep in flight):
convert the horizon-limited gain into a per-step contraction guarantee.
(ii) \emph{Conditioned-floor gap}: close the factor 2--3 between the trained
conditioned model and its information-theoretic ceiling (longer training,
richer $\shat$ features, angular harmonics for the $r_3$ remainder).
(iii) \emph{Weight-tied recursion cell}: one cell unrolled $p$ times
generalizes the estimator to any order $N^{(p)}$ with the exact $\partial_t$
recursion---no inner wall---at the price of a $\Delta^{-1}$ per order; the
exact-FFT vs learned-local $\Delta^{-1}$ decision is open, and is the same
decision as the deferred flux-form vs advective Jacobian consistency.
(iv) \emph{Decaying-turbulence members} (non-stationary: the $\taul$ identity
needs detrending) and obstacle flows.
(v) \emph{Generalization across schemes}: the entire construction used only
$(r,s)$ and the AB bracket; matching any linear multistep method to its
better-resolved self is mechanical from Prop.~\ref{prop:lte-ab4}'s template.

\section*{Reproducibility}
\begin{sloppypar}
The spectral QG solver follows \citet{sureshbabu2025}. Truncation operators:
\texttt{analysis/validate\_ab2cn2\_vs\_truth.py}; magnitudes and radii:
\texttt{analysis/measure\_truncation\_magnitudes.py},
\texttt{Theoretical\_guarantees/convergence\_radius.py}; microscale:
\texttt{check\_decorrelation\_time.py}; stencil walls:
\texttt{fd\_depth\_check.py}, \texttt{temporal\_fd\_floor\_deep.py}; error
propagation: \texttt{closure\_error\_propagation.py}; Wiener floors:
\texttt{check\_wiener\_floor.py}; stability regions:
\texttt{analysis/ab\_stability\_regions.py},
\texttt{analysis/vn\_stability.py},
\texttt{Theoretical\_guarantees/imex\_stability\_regions.py}; certificate:
\texttt{training/wiener\_certificate.py}; training:
\texttt{training/train\_deriv.py}, \texttt{training/train\_deriv\_rollout.py};
evaluation: \texttt{training/eval\_deriv\_by\_root.py},
\texttt{diagnostics/rollout\_aposteriori.py}.
\TODO{code/data availability statement + repository link.}
\end{sloppypar}

% ===================================================================== %
\appendix
\section{RK4 stage expansions}
\label{app:stages}
Throughout, $f=G(\omega^n)$, $J=G'(\omega^n)=\Lop+2B(\omega^n,\cdot)$, and the
exact identity \eqref{eq:frechet}: $G(\omega^n+p)=f+Jp+B(p,p)$, $B$ symmetric
bilinear.

\subsection{$k_2$}
$p=\tfrac h2f$:
$k_2=f+\tfrac h2Jf+\tfrac{h^2}4B(f,f)$.

\subsection{$k_3$}
$p=\tfrac h2k_2$. Linear part
$Jp=\tfrac h2Jf+\tfrac{h^2}4J^2f+\tfrac{h^3}8JB(f,f)$; quadratic part
$B(p,p)=\tfrac{h^2}4B(f,f)+\tfrac{h^3}4B(f,Jf)
+\tfrac{h^4}8B(f,B(f,f))+\tfrac{h^4}{16}B(Jf,Jf)+\bO(h^5)$. Hence
\begin{equation}
k_3=f+\tfrac h2Jf+\tfrac{h^2}4\big[J^2f+B(f,f)\big]
+\tfrac{h^3}8JB(f,f)+\tfrac{h^3}4B(f,Jf)
+\tfrac{h^4}8B(f,B(f,f))+\tfrac{h^4}{16}B(Jf,Jf).
\end{equation}

\subsection{$k_4$}
$p=hk_3$. Collecting linear and quadratic parts through $h^4$,
\begin{align}
k_4&=f+hJf+h^2\big[\tfrac12J^2f+B(f,f)\big]
+h^3\big[\tfrac14J^3f+\tfrac14JB(f,f)+B(f,Jf)\big]\notag\\
&\quad+h^4\big[\tfrac18J^2B(f,f)+\tfrac14JB(f,Jf)+\tfrac12B(f,J^2f)
+\tfrac12B(f,B(f,f))+\tfrac14B(Jf,Jf)\big].
\end{align}
No stage builds the depth-four chain, so the $h^4$ coefficient of $J^4f$ in
$k_4$ is zero; assembling $\tfrac h6(k_1+2k_2+2k_3+k_4)$ and subtracting from
the exact expansion gives Prop.~\ref{prop:lte-rk4}, per elementary
differential:
\begin{center}\small
\renewcommand{\arraystretch}{1.3}
\begin{tabular}{lccc}
\toprule
elementary differential & exact $\times h^5$ & RK4 $\times h^5$ &
$\tau_{\rm RK4}\times h^5$\\
\midrule
$J^4f$        & $1/120$ & $0$    & $1/120$\\
$J^2B(f,f)$   & $1/60$  & $1/48$ & $-1/240$\\
$JB(f,Jf)$    & $1/20$  & $1/24$ & $1/120$\\
$B(f,J^2f)$   & $1/15$  & $1/12$ & $-1/60$\\
$B(f,B(f,f))$ & $2/15$  & $1/8$  & $1/120$\\
$B(Jf,Jf)$    & $1/20$  & $1/16$ & $-1/80$\\
\bottomrule
\end{tabular}
\end{center}

% ===================================================================== %
\bibliographystyle{plainnat}
\begin{thebibliography}{99}

\bibitem[Ascher et~al.(1995)]{ascher1995}
U.~M. Ascher, S.~J. Ruuth, and B.~T.~R. Wetton.
Implicit--explicit methods for time-dependent partial differential equations.
\emph{SIAM J.\ Numer.\ Anal.}, 32(3):797--823, 1995.

\bibitem[Bar-Sinai et~al.(2019)]{barsinai2019}
Y.~Bar-Sinai, S.~Hoyer, J.~Hickey, and M.~P. Brenner.
Learning data-driven discretizations for partial differential equations.
\emph{Proc.\ Natl.\ Acad.\ Sci.}, 116(31):15344--15349, 2019.

\bibitem[Boyd(2001)]{boyd2001}
J.~P. Boyd. \emph{Chebyshev and Fourier Spectral Methods}. 2nd ed., Dover,
2001.

\bibitem[Butcher(2016)]{butcher2016}
J.~C. Butcher. \emph{Numerical Methods for Ordinary Differential Equations}.
3rd ed., Wiley, 2016.

\bibitem[Canuto et~al.(2006)]{canuto2006}
C.~Canuto, M.~Y. Hussaini, A.~Quarteroni, and T.~A. Zang.
\emph{Spectral Methods: Fundamentals in Single Domains}. Springer, 2006.

\bibitem[Charous \& Lermusiaux(2023)]{charous2023}
A.~Charous and P.~F.~J. Lermusiaux. Dynamically orthogonal Runge--Kutta
schemes with perturbative retractions for the dynamical low-rank
approximation. \emph{SIAM J.\ Sci.\ Comput.}, 45(2):A872--A897, 2023.

\bibitem[Frank et~al.(1997)]{frank1997}
J.~Frank, W.~Hundsdorfer, and J.~G. Verwer.
On the stability of implicit-explicit linear multistep methods.
\emph{Appl.\ Numer.\ Math.}, 25(2--3):193--205, 1997.

\bibitem[Guan et~al.(2022)]{guan2022}
Y.~Guan, A.~Chattopadhyay, A.~Subel, and P.~Hassanzadeh. Stable a posteriori
LES of 2D turbulence using convolutional neural networks. \emph{J.\ Comput.\
Phys.}, 458:111090, 2022.

\bibitem[Gupta \& Lermusiaux(2021)]{gupta2021}
A.~Gupta and P.~F.~J. Lermusiaux. Neural closure models for dynamical
systems. \emph{Proc.\ R.\ Soc.\ A}, 477:20201004, 2021.

\bibitem[Hairer et~al.(1993)]{hairer1993}
E.~Hairer, S.~P. N{\o}rsett, and G.~Wanner. \emph{Solving Ordinary
Differential Equations I: Nonstiff Problems}. 2nd ed., Springer, 1993.

\bibitem[Kochkov et~al.(2021)]{kochkov2021}
D.~Kochkov, J.~A. Smith, A.~Alieva, Q.~Wang, M.~P. Brenner, and S.~Hoyer.
Machine learning--accelerated computational fluid dynamics.
\emph{Proc.\ Natl.\ Acad.\ Sci.}, 118(21):e2101784118, 2021.

\bibitem[Lions et~al.(2001)]{lions2001}
J.-L. Lions, Y.~Maday, and G.~Turinici. R\'esolution d'EDP par un sch\'ema en
temps ``parar\'eel''. \emph{C.\ R.\ Acad.\ Sci.\ Paris}, 332:661--668, 2001.

\bibitem[Srinivasan et~al.(2024)]{srinivasan2024}
K.~Srinivasan, M.~D. Chekroun, and J.~C. McWilliams. Turbulence closure with
small, local neural networks: forced two-dimensional and $\beta$-plane flows.
\emph{J.\ Adv.\ Model.\ Earth Syst.}, 16:e2023MS003795, 2024.

\bibitem[Stetter(1978)]{stetter1978}
H.~J. Stetter. The defect correction principle and discretization methods.
\emph{Numer.\ Math.}, 29:425--443, 1978.

\bibitem[Suresh Babu et~al.(2025)]{sureshbabu2025}
A.~N. Suresh Babu, A.~Sadam, and P.~F.~J. Lermusiaux. \TODO{title}.
arXiv:2508.06678, 2025.

\bibitem[Vallis(2017)]{vallis2017}
G.~K. Vallis. \emph{Atmospheric and Oceanic Fluid Dynamics}. 2nd ed.,
Cambridge University Press, 2017.

\bibitem[Warming \& Hyett(1974)]{warming1974}
R.~F. Warming and B.~J. Hyett. The modified equation approach to the
stability and accuracy analysis of finite-difference methods. \emph{J.\
Comput.\ Phys.}, 14:159--179, 1974.

\bibitem[Wiener(1949)]{wiener1949}
N.~Wiener. \emph{Extrapolation, Interpolation, and Smoothing of Stationary
Time Series}. MIT Press, 1949.

\end{thebibliography}

\end{document}
